\section{Corrected dispersion}
\label{sec:dispersion}

The correction in this section is the cubic-frequency term in Poisson
summation. This is the dispersion mechanism used by Dunn and
Radziwi\l\l~\cite[Section 9, Propositions 9.1--9.2]{DR2024}.
Their squarefree majorant and cube-frequency correction guide the
argument. The noncube estimates and the Type~I mixed-term comparison
are proved here with the unconditional inputs of the preceding sections.
We keep track of the coefficient restrictions
and the frequency ranges needed here; the two exceptional ranges will be
handled by Proposition~\ref{prop:dual}.

Throughout the section, put
\begin{equation}
 AB\asymp X,\qquad X^{0.32}\le B\ll X^{1/2},\qquad
 H=\frac{B^2}{A},\qquad \mathcal Q=ABH^{1/3},\qquad L=\log X.
 \label{eq:dispersion-scales}
\end{equation}
All sums over letters denoting ideals use primary generators prime to
$3$, unless a full lattice is expressly specified. Let $\beta_b$ be
divisor-bounded, supported on squarefree $b$ with $Nb\asymp B$. When
$B>X^{0.40}$, suppose in addition that $\beta$ is a squarefree-restricted
convolution of a fixed number of prime sequences of the kind in
Proposition~\ref{prop:dual}, each prime having norm at least $X^c$ for
some fixed $c>0$. The weights on the individual prime dyads are smooth;
there is no additional cutoff coupling their products. Set
\begin{equation}
 \begin{aligned}
 \beta_b(u)&=\beta_b(Nb)^{iu},&
 S(u)&=\sum_b\beta_b(u)(Nb)^{-1/6},\\
 F_u(a)&=\sum_b\beta_b(u)G(b)\overline{\chi_b(a)}.
 \end{aligned}
 \label{eq:dispersion-polynomials}
\end{equation}
Let $W$ be a nonnegative radial smooth function, bounded in absolute
value and supported in a fixed annulus in $\mathbb C$, and write
$W_A(a)=W(a/\sqrt A)$.

The low-height assertion uses the following property of $\beta$:
\begin{quote}
\emph{(SW)} For every fixed conductor log-power and every fixed
norm-height log-power, the sum of $\beta_b$ against a nonprincipal cubic
Hecke character of the indicated conductor, with a norm twist of the
indicated height, is $O_D(BL^{-D})$ for every $D>0$. The same assertion
holds after imposing coprimality to any specified element of norm
$X^{O(1)}$.
\end{quote}
The implied constants are uniform in these characters, twists, and
exclusions. The fixed powers in the divisor bounds are chosen before any
roughness parameter below.

\begin{proposition}[Dispersion estimates]
\label{prop:dispersion}
Under the preceding hypotheses the following assertions hold.
\begin{enumerate}
\item If $H\gg1$ and $|u|\le X^{0.17}$, and the derivatives of $W$ are
bounded independently of $X$, then
\[
 \sum_a W_A(a)|F_u(a)|^2\ll \mathcal Q L^{O(1)}.
\]
\item Suppose $|u|\le L^{O(1)}$, $H\gg1$, and $\beta$ has property
\emph{(SW)} and is supported on elements all of whose prime factors have
norm greater than $L^{C_r}$. For every fixed $K>0$, if $C_r,C_g$ are
sufficiently large, then for $H\ge L^{C_g}$,
\begin{equation}
 \sum_a\mu^2(a)W_A(a)
 \left|F_u(a)-c_*\overline{G(a)}(Na)^{-1/6}S(u)\right|^2
 \ll \mathcal Q L^{-K}.
 \label{eq:dispersion-corrected}
\end{equation}
For $1\ll H<L^{C_g}$ one has the uncorrected estimate
\begin{equation}
 \sum_a W_A(a)|F_u(a)|^2
 \ll A\sum_b|\beta_b|^2+
 \mathcal Q\left\{\left(B^{-1}\sum_b|\beta_b|\right)^2+L^{-K}\right\}.
 \label{eq:dispersion-short}
\end{equation}
Here $1\ll H$ means only that $H$ is bounded below by a positive
constant. In these low-height assertions the derivatives of $W$ are
bounded independently of $X$.
\item Suppose $u=\pm t+u_0$, $T\le |t|\le2T$,
$2T+|u_0|\le X^{0.17}$, and
$|\partial^jW|\ll_j L^{Cj}$ for fixed $C$. For every fixed $K>0$, if
$T\ge L^{C_t}$ with $C_t$ sufficiently large, then
\begin{equation}
 \int_{T\le|t|\le2T}\sum_a W_A(a)|F_{\pm t+u_0}(a)|^2\frac{dt}{T}
 \le \int_{T\le|t|\le2T}\mathcal D_W\frac{dt}{T}
       +O(\mathcal Q L^{-K}),
 \label{eq:dispersion-height}
\end{equation}
where the exact diagonal is
\[
 \mathcal D_W=\sum_b|\beta_b|^2
                  \sum_{(a,b)=1}W_A(a).
\]
No lower bound on $H$ is required here. When $B\le X^{0.40}$ the
restriction on $u_0$ can be omitted.

There is also an absolute sufficiently small $\gamma>0$ such that, when
$B\le X^{0.40}$, $T\gg X^{0.01}$, and
$|\partial^jW|\ll_j X^{j\gamma}$, the corresponding estimate is
\begin{equation}
 \int_{T\le|t|\le2T}\sum_a W_A(a)|F_{\pm t+u_0}(a)|^2\frac{dt}{T}
 \le \int_{T\le|t|\le2T}\mathcal D_W\frac{dt}{T}
       +O(\mathcal Q X^{-10\gamma}).
 \label{eq:dispersion-power}
\end{equation}
This $\gamma$ is independent of the fixed powers in the divisor bounds.
In the last two estimates, a radial support of width $O(J^{-1})$ gives
\[
 \mathcal D_W\ll \frac AJ\sum_b|\beta_b|^2
\]
for $J\le L^{O(1)}$, or $J\le X^\gamma$, respectively. In particular,
when $B\le X^{0.40}$, the coefficients may be restricted sharply to any
norm cell without changing these assertions.
\end{enumerate}
\end{proposition}

\begin{proof}
Opening the second moment isolates its diagonal; Poisson summation
expresses the off-diagonal part in frequencies.
For \eqref{eq:dispersion-corrected}, the cube
contributions identified below produce the model square, while the
noncube contribution is small. The Type~I comparison supplies the
matching mixed term, giving cancellation in the corrected square.
For the height averages, averaging makes the cube contributions small
as well, leaving the exact diagonal needed for norm localization.
We establish the shared frequency bounds before evaluating the cube
main term and completing that cancellation.

All arbitrarily small power losses below come from Heath-Brown's cubic large
sieve~\cite{HB2000}, in the forms of
Theorem~\ref{thm:cls} and Lemma~\ref{lem:cubic-extensions}; they will be
chosen smaller than the fixed power gaps. For logarithmic savings at
small conductors we instead use Lemma~\ref{lem:ordinary-sieve}, with no
such loss.

\paragraph{Squarefree majorization.}
Only for \eqref{eq:dispersion-corrected}, majorize the squarefree weight
in the positive $|F_u(a)|^2$ term by
\begin{equation}
 \mu^2(a)\le
 \left(\sum_{\substack{c^2\mid a\\Nc\le L^{D_0}}}\mu(c)\right)^2
 =\sum_{d^2\mid a}\lambda_d.
 \label{eq:dispersion-sieve}
\end{equation}
The inequality is equality on squarefree $a$. Expanding the square shows
that $d$ is squarefree, $Nd\le L^{2D_0}$,
$|\lambda_d|\le\tau(d)^2$, and $\lambda_d=\mu(d)$ for $Nd\le L^{D_0}$.
For the uncorrected estimates take only $d=1$, with coefficient one.
The mixed and model terms of the corrected square will retain their
true squarefree weights.

The diagonal is $\mathcal D_W$ when $d=1$. In the sieved case it is
$O(AB L^{O(1)})$: for example, sum
$|\lambda_d|A/(Nd)^2$ and use the coefficient second moment. This is
$O(\mathcal Q L^{-K})$ once $H\ge L^{C_g}$ with $C_g$ large enough.

\paragraph{Poisson summation.}
For an off-diagonal pair write
\[
 f=(b_1,b_2),\qquad b_i=fp_i,\qquad F=Nf,
\]
so that $p_1,p_2$ are coprime, squarefree, prime to $f$, and have norm
$\asymp B/F$. The summand vanishes unless $(b_1b_2,d)=1$. Remove the
remaining restriction $(a,f)=1$ by Mobius inversion, and write
\[
 a=d^2mx,\qquad m\mid f,\qquad M=Nm,\qquad D_d=Nd.
\]
Use normalized plane measure and Fourier transform
\[
 d\mathfrak m(z)=\frac2{\sqrt3}\,d\Re z\,d\Im z,
 \qquad
 \widehat W(y)=\int_{\mathbb C}W(z)e(-\operatorname{Tr}(zy))
                                      \,d\mathfrak m(z).
\]
Poisson summation gives, for this pair and these $d,m$, the following
expression, to be multiplied by $\lambda_d\mu(m)$:
\begin{equation}
\begin{split}
 &\frac{A}{9D_d^2M}
  \frac{\beta_{fp_1}(u)\overline{\beta_{fp_2}(u)}}{|p_1p_2|}
  \sum_{\substack{h\in\mathcal O\\h\ne0}}
  e\left(\operatorname{Tr}\frac{h}{3\lambda}\right)
  \widehat W\left(
    \frac{h\sqrt A}{3\lambda D_d\sqrt M\,p_1p_2}\right)\\
 &\hspace{35mm}\cdot
  \chi_{p_1}\bigl(dh(fm)^{-1}\bigr)
  \overline{\chi_{p_2}\bigl(dh(fm)^{-1}\bigr)}.
\end{split}
\label{eq:poisson}
\end{equation}
The inverses here and below are residue-class notation.

Here are the arithmetic details of this identity. The trace-dual lattice
of $\mathcal O$ is $\lambda^{-1}\mathcal O$. Thus the modulus
$3p_1p_2$ has index $9N(p_1p_2)$, and the dilation contributes
$A/(D_d^2M)$. Since $d,m$ are primary, $x$ must be primary. Put
$q=p_1p_2$ and represent this coset by $x=q+3s$, $s$ modulo $q$.
Its additive factor is
\[
 e\left(\operatorname{Tr}\frac{h}{3\lambda}\right)
 e\left(\operatorname{Tr}\frac{sh}{\lambda q}\right).
\]
The local Gauss sums have magnitude $|q|$. Their normalized factors
cancel those in $G(fp_1)\overline{G(fp_2)}$, by \eqref{eq:mult}.
The CRT factors at $p_1$ and $p_2$ cancel by reciprocity, and the factor
$3\lambda=(-\lambda)^3$ is a cube. What remains from $f$, $m$, and
$d^2$ is exactly
$\chi_{p_1}(dh/(fm))\overline{\chi_{p_2}(dh/(fm))}$, since
$\overline{\chi(d^2)}=\chi(d)$. The common local factors at $f$ were
the indicator $(a,f)=1$, explaining its Mobius removal. At frequency
zero at least one nontrivial character modulo $p_i$ has zero sum,
because the original pair was unequal. Finally, radiality permits the
complex dilation by $d^2m$ to be replaced by its absolute value.

\paragraph{Discarding large common divisors.}
Fix a sufficiently small positive $\delta$, to be specified after the
frequency estimates, and first remove $F\ge X^\delta$. The natural
dual norm scale in \eqref{eq:poisson} is
\begin{equation}
 H'=\frac{HD_d^2M}{F^2}.
 \label{eq:dispersion-dual-scale}
\end{equation}
If $H<X^{\delta/2}$, then $M\le F$ gives
$H'\le X^{-\delta/2+o(1)}$. Every nonzero frequency is in a rapidly
decreasing Fourier tail, so the contribution is power-negligible. This
remains so for the power-width weights after choosing $\gamma$
sufficiently small in terms of $\delta$.

If $H\ge X^{\delta/2}$, invert the function
$1_{N(b_1,b_2)\ge X^\delta}$ on the divisor lattice. Its divisor
coefficients are bounded by $\tau(e)$ and vanish for $Ne<X^\delta$.
For each fixed $e$, the resulting full form is bounded by a sum of
squares with inner length $B/Ne$: the common factor
$\chi_e(a)$ drops out in absolute value, and \eqref{eq:mult} merely
rephases the coefficients. The sieve weight costs at most a small power.
By \eqref{eq:outer-extension}, a dyad $Ne\asymp E$ therefore costs
\[
 X^\epsilon E\left(A+(AB/E)^{2/3}\right)\frac BE
 =X^\epsilon\left(AB+\mathcal Q E^{-2/3}\right).
\]
Here $A\gg B/E$ allows the last term of
\eqref{eq:outer-extension} to be absorbed. Relative to $\mathcal Q$,
the first term saves $H^{-1/3}\le X^{-\delta/6}$ and the second saves
$E^{-2/3}\le X^{-2\delta/3}$. Removing a diagonal costs at most
$X^\epsilon AB$, with the same saving. Summing the logarithmically many
dyads proves the required power saving.

\paragraph{Separating the remaining pairs.}
For $F<X^\delta$, both $p_i$ are nontrivial. In
\eqref{eq:poisson} remove their mutual coprimality by
$\sum_{l\mid(p_1,p_2)}\mu(l)$, and set
\[
 p_i=ln_i,\qquad V_l=Nl,\qquad B'=B/(FV_l).
\]
For clarity, we extend the \emph{displayed expression} in
\eqref{eq:poisson} to all pairs before this Mobius inversion; we do not
assert the original Poisson identity for artificial noncoprime pairs.
The inversion returns precisely the original coprime pairs. Squarefree
support forces $(l,n_i)=1$. The common character factors then multiply
to $|\chi_l(dh/(fm))|^2$, an exclusion on $h$ alone. All other fixed
coprimality conditions may be put into the individual coefficients.
Artificial equal pairs cancel in the exact Mobius identity; the only
exceptional pair $p_1=p_2=1$ is impossible in this range.

On a frequency dyad $Nh\asymp J$, put $k=J/H'$. Mellin separation of
the radial transform and the norm denominators, followed by Cauchy in
$h$, reduces the estimate to shifted squares
\begin{equation}
 \sum_{\substack{h\in\mathcal O\\Nh\asymp J}}
 \left|\sum_n\beta_{fln}(Nn)^{iu'}
             \chi_n\bigl(dh(fm)^{-1}\bigr)\right|^2,
 \label{eq:dispersion-separated}
\end{equation}
with fixed exclusions understood. The shift $u'$ is $u$ plus a Mellin
variable. A smooth cutoff equal to one on all possible scaled norm
ratios lets us place every additional smooth factor into the separated
kernel. Thus at $f=l=1$ no new cutoff on the prime products is imposed:
both the prime-convolution hypothesis and (SW) are preserved.

We record the transform cost explicitly. On the compact sets of positive
scaled norms, integration by parts in the plane Fourier transform and
then in the two log-norm variables gives, for arbitrary fixed $R,M_0$,
a separated-kernel bound of the form
\[
 C_{R,M_0}(W)(1+k)^{-R}
 (1+|v_1|+|v_2|)^{-M_0}.
\]
Here $C_{R,M_0}(W)$ is bounded by finitely many derivative seminorms of
$W$. It is a fixed log-power for log-scale weights and $X^{O(\gamma)}$
for power-width weights. We can truncate $k$ and all shifts at
$X^{\delta'}$, for any sufficiently small fixed $\delta'>0$, with
power-negligible tails, choosing finitely many integrations by parts
first and $\gamma$ afterwards. Inside the truncation we retain these
decaying factors, so their integrals, even with fixed polynomial
weights in $k$ and the shifts, cost only log-powers or $X^{O(\gamma)}$.
All shifts are independent of the averaging variable $t$. Different
shifts on the two sides of Cauchy present no problem. There are $O(L)$
frequency dyads.

Use $J^{1/3}B'^2$ as a benchmark for
\eqref{eq:dispersion-separated}. The prefactor in
\eqref{eq:poisson} has size $AF/(D_d^2MB)$. Dividing its product with
this benchmark by $\mathcal Q$ gives
\begin{equation}
 \frac{AF}{D_d^2MB}\frac{J^{1/3}B'^2}{\mathcal Q}
 =k^{1/3}D_d^{-4/3}M^{-2/3}F^{-5/3}V_l^{-2}.
 \label{eq:dispersion-benchmark}
\end{equation}
The $f$ sum with $m\mid f$, and the $d$ sum with its divisor-bounded
sieve coefficients, converge with these powers.

First discard $V_l\ge X^\delta$.
Equation~\eqref{eq:outer-extension} bounds the square by
\[
 X^\epsilon\bigl(J+(JB')^{2/3}+J^{1/3}B'\bigr)B'.
\]
Its ratios to the benchmark are $J^{2/3}/B'$, $(J/B')^{1/3}$, and $1$.
On the retained transforms,
$J\ll X^{2\delta'}B/F$, since $H\ll B$, $M\le F$, and $D_d$ is
polylogarithmic. Also $B/F\gg X^{0.3}$. After summing the $l$ in a
dyad, the respective costs from the $V_l^{-2}$ factor in
\eqref{eq:dispersion-benchmark} are bounded, up to small powers, by
\[
 X^{4\delta'/3}(B/F)^{-1/3},\qquad
 X^{2\delta'/3}V_l^{-2/3},\qquad V_l^{-1}.
\]
Choose $\delta'$ sufficiently small relative to $\delta$. Each term
then saves a fixed power, including after all divisor and dyadic sums.

\paragraph{The large-conductor frequency ranges.}
For the remaining $l$, factor a nonzero frequency as
\[
 h=\zeta\lambda^i p q^2 j^3,\qquad i\in\{0,1,2\},
\]
where $p,q$ are coprime squarefree primary elements. A fixed convention
on units makes the multiplicity bounded; the cube part may contain
powers of $\lambda$. In dyads put $Np\asymp P$, $Nq\asymp Q$.
Then $PQ^2\ll J$ and $Nj\asymp(J/(PQ^2))^{1/3}$, up to bounded
factors. For each fixed $j$ the relevant benchmark is
\begin{equation}
 B'^2(PQ^2)^{1/3}.
 \label{eq:dispersion-pq-benchmark}
\end{equation}
The number of $j$ restores the benchmark $J^{1/3}B'^2$. The cube part
only excludes $(n,j)>1$; the other fixed symbols are coefficient
twists. The inverse twist can, if desired, be written using $(fm)^2$
on its coprime support.

We claim a fixed power saving over
\eqref{eq:dispersion-pq-benchmark} whenever $PQ^2\ge X^{0.003}$.
Here is the full limiting-exponent check. Let $\delta,\delta'$ shrink
to zero along a hypothetical sequence without a power saving. Then
$B'=B^{1+o(1)}$ and $PQ^2\le B^{1+o(1)}$. If $P$ has larger limiting
exponent than $Q$, fix $q$ and use Theorem~\ref{thm:cls} on $p$.
The bound $QB(P+B+(PB)^{2/3})$, divided by
$B^2(PQ^2)^{1/3}$, has the three ratios
\begin{equation}
 \frac{P^{2/3}Q^{1/3}}B,\qquad
 (Q/P)^{1/3},\qquad (PQ/B)^{1/3}.
 \label{eq:dispersion-p-dominant}
\end{equation}
Under $PQ^2\le B^{1+o(1)}$, all save unless
$P=B^{1+o(1)}$, $Q=B^{o(1)}$. If $Q$ has larger exponent, reverse
the roles, conjugating the cubic character as necessary. The ratios are
\begin{equation}
 \frac{P^{2/3}Q^{1/3}}B,\qquad
 (P/Q)^{2/3},\qquad (P^2/B)^{1/3},
 \label{eq:dispersion-q-dominant}
\end{equation}
all with strict gaps. If the exponents agree, write
$P=Q=B^{v+o(1)}$. Lemma~\ref{lem:ordinary-sieve} gives the bound
$(B+(PQ)^2)B$, whose ratio is
\[
 B^{-v+o(1)}+B^{3v-1+o(1)}.
\]
Since $0<v\le1/3$, this saves unless $v=1/3$; the endpoint $v=0$
is ruled out by $PQ^2\ge X^{0.003}$.

The two exceptional configurations are exactly those of
Proposition~\ref{prop:dual}. Both force $PQ^2=B^{1+o(1)}$, hence
$J\ge B^{1-o(1)}$, $H\ge B^{1-o(1)}$, and
$B=X^{1/2+o(1)}$. Thus the required prime-convolution hypothesis is
available. Since every supported prime has norm at least $X^c$, choosing
$\delta<c$ forces $f=l=1$ eventually in this limiting test.
The remaining $d$, cube part $j$, and fixed twists or exclusions have
subpower norm. Moreover
\[
 1+|u'|\le X^{0.17}+X^{\delta'}+1\le B^{0.36}
\]
eventually. Proposition~\ref{prop:dual} therefore supplies the missing
strict power saving. To check the normalization of this saving, write
$\delta_D,\eta_D$ for the constants in that proposition and choose an
exceptional exponent radius $h_0<\min(\delta_D/4,\eta_D/4)$.
In either such neighborhood,
$B^2(PQ^2)^{1/3}\ge B^{7/3-h_0}$.
Thus the bound $B^{7/3-\eta_D}$ still saves at least
$B^{3\eta_D/4}$ relative to the local benchmark; the allowed extra
twists and exclusions are restricted to the same smaller neighborhood.
Its fixed-neighborhood interpretation and
compactness of the exponent region yield one positive saving for some
fixed sufficiently small $\delta,\delta'$. The arbitrary small-power
losses are chosen after this gap.

If $B\le X^{0.40}$, the two exceptional configurations cannot occur:
$H=B^3/X\le X^{0.20}$, whereas $B\ge X^{0.32}$, so $H$ cannot be
$B^{1-o(1)}$. Consequently the gaps in this range are absolute. We may
choose $\delta,\delta'$, and subsequently $\gamma$, separately for
this range, independently of the prime-convolution parameters or of
the fixed divisor exponents. No restriction on the norm twists is
needed in this case.

\paragraph{Small conductors and averaging in height.}
Suppose now that $PQ^2<X^{0.003}$. The primitive characters indexed by
the coprime squarefree pair $(p,q)$ have conductor away from $3$ equal
to $pq$, and the two exponents distinguish them. All fixed symbols and
exclusions can be put into the coefficients. Since
$(PQ)^2\le X^{0.006}\ll B'$, Lemma~\ref{lem:ordinary-sieve} bounds
the ratio to \eqref{eq:dispersion-pq-benchmark} by
\begin{equation}
 L^{O(1)}\tau(fl)^{O(1)}(PQ^2)^{-1/3}.
 \label{eq:dispersion-small-conductor}
\end{equation}
The divisor factor is summable against
\eqref{eq:dispersion-benchmark}. This gives any desired log saving
once $PQ^2$ exceeds a sufficiently large fixed log-power. Without any
improvement for the smaller $P,Q$, it gives a log-power bound throughout.
Together with the preceding power-saving ranges and the diagonal, this
proves the coarse assertion when $H\gg1$.

For the height average, fix one character row and collect equal norms
in its inner polynomial. Divisor moments bound the squared coefficient
sum by $B'L^{O(1)}\tau(fl)^{O(1)}$. The ordinary mean-value inequality
therefore gives
\[
 \int_{T\le|t|\le2T}|\text{inner polynomial}|^2\frac{dt}{T}
 \ll (B'+B'^2/T)L^{O(1)}\tau(fl)^{O(1)}.
\]
After summing the $O(PQ)$ rows, the ratio to
\eqref{eq:dispersion-pq-benchmark} is
\begin{equation}
 L^{O(1)}\tau(fl)^{O(1)}P^{2/3}Q^{1/3}
       \left(T^{-1}+B'^{-1}\right).
 \label{eq:dispersion-height-small}
\end{equation}
Use \eqref{eq:dispersion-small-conductor} above a sufficiently large
log-power of $PQ^2$, and \eqref{eq:dispersion-height-small} below it.
Taking $C_t$ large absorbs all fixed logarithmic output and transform
costs, proving \eqref{eq:dispersion-height}.

For \eqref{eq:dispersion-power}, throughout this small-conductor range
\[
 P^{2/3}Q^{1/3}\le(PQ^2)^{2/3}<X^{0.002},\qquad
 T^{-1}\ll X^{-0.01},\qquad B'^{-1}\ll X^{-0.3}.
\]
There is thus an absolute power gap, at least $0.008$ before arbitrarily
small losses, in \eqref{eq:dispersion-height-small}. The other frequency
ranges already have absolute gaps when $B\le X^{0.40}$. Choose
$\gamma$ small enough that these gaps absorb all the finitely many
transform seminorm costs and the requested factor $X^{10\gamma}$.
Fixed divisor powers are harmless because their losses can be taken
arbitrarily small. This proves \eqref{eq:dispersion-power}.

Neither argument used $H\gg1$: if $H'$ is small, the nonzero frequency
dyads merely have large $k=J/H'$ and are in the Fourier tail. In
particular, no cube lattice asymptotic is used for the height averages.
For $B\le X^{0.40}$ only coefficient moments and large sieves were
used, which proves the asserted freedom of cell restrictions and of
$u_0$. Finally an annulus of radial width $O(J^{-1})$ contains
$O(A/J+\sqrt A+1)$ lattice points. Since $A\gg X^{1/2}$ and $\gamma$
is small, this is $O(A/J)$ in the stated ranges, proving the localized
diagonal bound.

\paragraph{Low-height nonprincipal frequencies.}
We now assume roughness and (SW). If $f$ or $l$ is nonunit, its norm
exceeds $L^{C_r}$. The convergent tails in
\eqref{eq:dispersion-benchmark}, with the fixed divisor factors, give
$O(\mathcal Q L^{-K'})$ for every desired $K'$ by taking $C_r$ large.
For example, the $f$ tail decays as $F^{-2/3}$ and the $l$ tail as
$V_l^{-1}$ on dyads, up to fixed divisor powers. The transform and
coefficient complexity costs do not grow with $C_r$.

It remains to consider $f=l=m=1$. Exclude precisely the frequencies
for which $dh$ is a cube. If $PQ^2$ is bounded by a fixed log-power,
the remaining character $n\mapsto\chi_n(dh)$ is nonprincipal, with
primitive conductor bounded by a fixed log-power. The cube part only
adds the exclusion $(n,j)=1$, permitted in (SW). Hence (SW) gives any
required log saving in \eqref{eq:dispersion-pq-benchmark}, taking its
exponent large enough to absorb the number of rows and the other
logarithmic costs. For Mellin shifts beyond a sufficiently large
log-power use rapid transform decay and the coarse bounds already
proved for the separated squares. Together with
\eqref{eq:dispersion-small-conductor} and the large-conductor power
savings, this leaves only the cube frequencies, with error
$O(\mathcal Q L^{-K'})$.

\paragraph{The cube frequencies and their constant.}
Since $d$ is squarefree, $dh$ is a cube exactly when
\[
 h=d^2j^3,\qquad j\in\mathcal O\setminus\{0\}.
\]
Every such $h$ has exactly three preimages, from the three cube roots
of unity. The additive phase in \eqref{eq:poisson} is one: a cube in
$\mathcal O$ is congruent modulo $3$ to a rational integer, and
$d\equiv1\pmod3$. Choose $C_r>2D_0$, so $d$ is coprime to every
supported $b$.

Temporarily discard $(j,b_1b_2)=1$. Substituting $h=d^2j^3$ into
\eqref{eq:poisson} cancels the magnitude of $d^2$ against $D_d$ in
the Fourier argument. A rotation is immaterial by radiality. The cube
frequency sum, counting each frequency once, is therefore
\begin{equation}
 \frac13\left(\frac{27N(b_1b_2)}A\right)^{1/3}
       \int_{\mathbb C}\widehat W(z^3)\,d\mathfrak m(z)
       +O(H^{1/6}+1).
 \label{eq:dispersion-cube-lattice}
\end{equation}
Indeed the lattice radius is comparable to $H^{1/6}$, and
$z\mapsto\widehat W(z^3)$ is Schwartz. The elementary lattice
Riemann-sum error is $O(H^{1/6}+1)$; deleting the origin costs $O(1)$.
For $H\ge L^{C_g}$ this is a relative $O(H^{-1/6})$ saving.
For any prime $\pi$, rapid decay bounds the nonzero multiples of $\pi$
in this lattice sum by $O(H^{1/3}/N\pi)$, also when the lattice
spacing exceeds the radius. Restoring coprimality thus costs
\[
 O\left(H^{1/3}\sum_{\pi\mid b_1b_2}\frac1{N\pi}\right),
\]
which saves arbitrarily many log-powers by roughness. For
\eqref{eq:dispersion-short}, with $d=1$, only the absolute bound
$O(H^{1/3})$ for this sum is needed. Its contribution is directly
\[
 \ll \frac AB H^{1/3}\left(\sum_b|\beta_b|\right)^2
 =\mathcal Q\left(B^{-1}\sum_b|\beta_b|\right)^2.
\]
Together with the diagonal and the already bounded noncube terms, this
proves \eqref{eq:dispersion-short}.

We evaluate the integral in \eqref{eq:dispersion-cube-lattice} exactly.
The cubing map has Jacobian $9|z|^4$ and generic multiplicity three, so
\[
 \int\widehat W(z^3)\,d\mathfrak m(z)
 =\frac13\int\widehat W(y)|y|^{-4/3}\,d\mathfrak m(y).
\]
For the trace Fourier pairing its Riesz-kernel constant, including
$1/3$, is
\[
 \frac13\frac2{\sqrt3}\pi(2\pi)^{-2/3}
       \frac{\Gamma(1/3)}{\Gamma(2/3)}
 =\frac{(2\pi)^{4/3}}{9\Gamma(2/3)^2}=c_*^2.
\]
For completeness, represent $|y|^{-4/3}$ as a Gamma integral of
Gaussians and apply the two-dimensional Gaussian Fourier integral;
the pairing $\operatorname{Tr}(zy)=2\Re(zy)$ replaces the ordinary
Fourier frequency by twice a reflected vector, while
$d\mathfrak m=(2/\sqrt3)\,d\Re y\,d\Im y$ supplies the other factor.
The last equality uses
$\Gamma(1/3)\Gamma(2/3)=2\pi/\sqrt3$. Gaussian regularization
justifies the interchange, since the singularity at zero is locally
integrable and $\widehat W$ is rapidly decreasing. We conclude that
\begin{equation}
 \int\widehat W(z^3)\,d\mathfrak m(z)
 =c_*^2\int W(y)|y|^{-2/3}\,d\mathfrak m(y).
 \label{eq:dispersion-riesz}
\end{equation}

Combining \eqref{eq:dispersion-cube-lattice} with the prefactor in
\eqref{eq:poisson} gives the scaled factor
\[
 \frac{A^{2/3}}{9D_d^2}(Nb_1Nb_2)^{-1/6}.
\]
Moreover,
\[
 \sum_d\frac{\lambda_d}{(Nd)^2}
 =\sum_d\frac{\mu(d)}{(Nd)^2}+O(L^{-D_0/2}),
\]
by $\lambda_d=\mu(d)$ up to $L^{D_0}$ and absolute convergence with
the divisor bound. The primary coset has density $1/9$. Ordinary
squarefree lattice counting, by square-divisor inversion, therefore
identifies the total cube term as
\begin{equation}
 c_*^2|S(u)|^2\sum_a\mu^2(a)W_A(a)(Na)^{-1/3}
       +O(\mathcal Q L^{-K'}),
 \label{eq:dispersion-cube-main}
\end{equation}
on taking $D_0,C_r,C_g$ sufficiently large. The counting error is a
power saving: a smooth annular squarefree count has error
$O_\epsilon(A^{1/2+\epsilon})$, and the extra weight contributes
$A^{-1/3}$ instead of the main scale $A^{2/3}$.

\paragraph{Cancellation in the corrected square.}
Open the square in \eqref{eq:dispersion-corrected}. The preceding work
bounds its positive $|F_u|^2$ part by the main term in
\eqref{eq:dispersion-cube-main} and $O(\mathcal Q L^{-K'})$.
For its mixed term retain the true squarefree weight and use
\eqref{eq:mult}:
\[
 \sum_a\mu^2(a)W_A(a)(Na)^{-1/6}G(a)F_u(a)
 =\sum_{a,b}\beta_b(u)W_A(a)(Na)^{-1/6}G(ab).
\]
Apply Proposition~\ref{prop:typeI} at $R=B$, $U=A$, which is allowed
since $B\ll X^{1/2}<X^{0.51}$. It replaces this by
\[
 c_*\sum_{a,b}\beta_b(u)\mu^2(ab)
                  W_A(a)(Na)^{-1/3}(Nb)^{-1/6},
\]
with error $O(A^{-1/6}X^{5/6-\eta})$. Since
$|S(u)|\ll B^{5/6}L^{O(1)}$, multiplication by $\overline{S(u)}$
leaves error
\[
 A^{-1/6}B^{5/6}X^{5/6-\eta}L^{O(1)}
 =\mathcal Q X^{-\eta}L^{O(1)}.
\]
In the model sum, remove the mutual coprimality of $a,b$ at logarithmic
cost. A shared prime has norm greater than $L^{C_r}$, and the unweighted
pair count saves a factor bounded by
$\sum_{N\pi>L^{C_r}}(N\pi)^{-2}$; coefficient moments and Cauchy
retain an arbitrarily large log saving as $C_r$ grows. Thus the mixed
main term, after multiplication by $c_*\overline{S(u)}$, is exactly
the main term in \eqref{eq:dispersion-cube-main}, within the desired
error. The model square has that same main term because
$|G(a)|^2=1$ on squarefree $a$. Its coefficients in the expanded
corrected square are consequently $1,-2,1$, which cancel.

Choose $D_0$ first, then $C_r$ large enough for roughness and the sieve
coprimality, and $C_g$ large enough for the diagonal and lattice errors;
use (SW) with the finitely many required conductor, twist, and saving
powers. This proves \eqref{eq:dispersion-corrected} and completes the
proof of all assertions.
\end{proof}

\begin{remark}[The logarithmic transition]
\label{rem:dispersion-transition}
The uncorrected estimate retains the useful prime sparsity. If $b$ is
one prime and $a$ is a product of two primes, each in its specified dyad
of length $X^{1/3}L^{O(1)}$, bounded weights give a bilinear bound
$O(X^{5/6}L^{-3/2})$ for $H\gg1$. Proposition~\ref{prop:bilinear-low}
will derive this from \eqref{eq:dispersion-short}, retaining the
different prime-density factors on the two sides. This half-power
of logarithmic saving beyond $X^{5/6}/L$ is what allows the
three-prime transition to survive a fixed power of $\log L$ in the
number of pieces. No power saving at that transition is required.
\end{remark}
